\documentclass[10pt,a4paper]{article}
\usepackage[a4paper,margin=14mm,top=12mm,bottom=14mm]{geometry}
\usepackage{fontspec}
\setmainfont{Liberation Sans}
\usepackage{polyglossia}
\setdefaultlanguage{french}
\usepackage{amsmath}
\usepackage{xcolor}
\usepackage{graphicx}
\usepackage[most]{tcolorbox}
\usepackage{tabularx,booktabs,array,multirow}
\usepackage{enumitem}
\usepackage{fancyhdr}
\usepackage{lastpage}
\usepackage[hidelinks]{hyperref}
\definecolor{uimm}{HTML}{005677}
\definecolor{cyan}{HTML}{13A2DD}
\definecolor{fond}{HTML}{E6F2FA}
\definecolor{ok}{HTML}{2E7D32}
\definecolor{ko}{HTML}{B3261E}
\definecolor{violet}{HTML}{7B1FA2}
\setlength{\parindent}{0pt}
\setlength{\parskip}{3pt}
\setlist{nosep}
\setlist[itemize]{leftmargin=13pt,label=\textcolor{cyan}{\textbullet}}
\newcommand{\titre}[1]{\vspace{5pt}{\color{uimm}\large\bfseries #1}\par\vspace{-3pt}{\color{cyan}\rule{\linewidth}{1.2pt}}\par\vspace{0pt}}
\tcbset{boite/.style={colback=fond,colframe=uimm,boxrule=0.6pt,arc=2pt,left=6pt,right=6pt,top=4pt,bottom=4pt},
        defi/.style={colback=violet!6,colframe=violet,boxrule=0.8pt,arc=2pt,left=6pt,right=6pt,top=4pt,bottom=4pt,
                     title={\bfseries Défi},coltitle=white,colbacktitle=violet,fonttitle=\small}}
\newcommand{\bandeau}[3]{%
\begin{tcolorbox}[colback=uimm,colframe=uimm,arc=2pt,left=4pt,right=8pt,top=4pt,bottom=4pt]
\begin{minipage}[c]{18mm}\setlength{\fboxsep}{1.2pt}\fcolorbox{white}{white}{\includegraphics[height=16.5mm]{logo_uimm.jpg}}\end{minipage}\hspace{4mm}%
\begin{minipage}[c]{\dimexpr\linewidth-22mm\relax}\color{white}
{\small #1}\\[1pt]{\LARGE\bfseries #2}\\[2pt]{\large #3}\hfill{\small CQPM Technicien qualité}
\end{minipage}
\end{tcolorbox}}
\newcommand{\reponse}[1][1]{\par\vspace{2pt}{\color{gray!60}\foreach \i in {1,...,#1}{\rule{\linewidth}{0.3pt}\par\vspace{9pt}}}}
\usepackage{pgffor}
\pagestyle{fancy}\fancyhf{}
\renewcommand{\headrulewidth}{0pt}
\newcommand{\pied}[1]{\fancyfoot[L]{\footnotesize\color{uimm}#1}\fancyfoot[R]{\footnotesize\color{uimm}Page \thepage{} sur \pageref{LastPage}}}

\usepackage{tikz}
\usepackage{longtable}
\renewcommand{\bandeau}[3]{%
\begin{tcolorbox}[colback=uimm,colframe=uimm,arc=2pt,left=4pt,right=8pt,top=4pt,bottom=4pt]
\begin{minipage}[c]{18mm}\setlength{\fboxsep}{1.2pt}\fcolorbox{white}{white}{\includegraphics[height=16.5mm]{logo_uimm.jpg}}\end{minipage}\hspace{4mm}%
\begin{minipage}[c]{\dimexpr\linewidth-22mm\relax}\color{white}
{\small #1}\\[1pt]{\LARGE\bfseries #2}\\[2pt]{\large #3}
\end{minipage}
\end{tcolorbox}}
% en-tête d'une fiche : code, titre, question, bloc, durée, organisation, niveau
\newcommand{\entete}[7]{%
\begin{tcolorbox}[colback=uimm,colframe=uimm,arc=2pt,left=4pt,right=6pt,top=4pt,bottom=4pt]
\begin{minipage}[c]{18mm}\setlength{\fboxsep}{1.2pt}\fcolorbox{white}{white}{\includegraphics[height=16.5mm]{logo_uimm.jpg}}\end{minipage}\hspace{3mm}%
\begin{minipage}[c]{\dimexpr\linewidth-58mm\relax}\color{white}\raggedright
{\small Base SPC · #4}\\[1pt]{\Large\bfseries #1 · #2}\\[2pt]{\itshape #3}
\end{minipage}\hfill
\begin{minipage}[c]{33mm}\color{white}\small\raggedleft
Durée : \textbf{#5}\\ \textbf{#6}\\ Niveau #7 sur 3
\end{minipage}
\end{tcolorbox}}
\newcommand{\tagE}{{\setlength{\fboxsep}{1.5pt}\colorbox{cyan}{\color{white}\scriptsize\bfseries Ellistat}}}
\newcommand{\tagW}{{\setlength{\fboxsep}{1.5pt}\colorbox{uimm}{\color{white}\scriptsize\bfseries Page web}}}
\tcbset{pred/.style={colback=orange!6,colframe=orange!80!black,boxrule=0.6pt,arc=2pt,left=6pt,right=6pt,top=3pt,bottom=1pt,
                    title={\bfseries Je prédis, avant de commencer},coltitle=white,colbacktitle=orange!80!black,fonttitle=\small},
        retenir/.style={colback=fond,colframe=uimm,boxrule=0.6pt,arc=2pt,left=6pt,right=6pt,top=3pt,bottom=3pt,
                    title={\bfseries À retenir},coltitle=white,colbacktitle=uimm,fonttitle=\small},
        plus/.style={colback=violet!6,colframe=violet,boxrule=0.6pt,arc=2pt,left=6pt,right=6pt,top=3pt,bottom=3pt,
                    title={\bfseries Pour aller plus loin},coltitle=white,colbacktitle=violet,fonttitle=\small},
        act/.style={colback=uimm!10,colframe=uimm,boxrule=0.5pt,arc=2pt,left=5pt,right=5pt,top=2pt,bottom=2pt,before skip=7pt,after skip=3pt},
        verif/.style={colback=orange!6,colframe=orange!80!black,boxrule=0.6pt,arc=2pt,left=6pt,right=6pt,top=3pt,bottom=3pt,
                    title={\bfseries À vérifier sur poste avant la première séance},coltitle=white,colbacktitle=orange!80!black,fonttitle=\small}}
\newcommand{\fichier}[1]{\texttt{\small #1}}
\newcommand{\menu}[1]{\textsf{\textbf{#1}}}
\newcommand{\champ}[1]{\rule{#1}{0.3pt}}

\pied{Base SPC · fiche G1 · Le pied à coulisse sur 150 mm}
\begin{document}
\thispagestyle{fancy}
\entete{G1}{Le pied à coulisse sur 150 mm}{L'erreur du pied à coulisse est-elle la même à 10 mm et à 150 mm ?}{Mallettes GUNT : plaques et arbres}{45 min}{en binôme}{2}
\begin{tcolorbox}[boite]\textbf{Le but.} Mesurer sept longueurs étalons de 10 à 150 mm au pied à coulisse, voir si l'erreur grandit avec la longueur, puis la comparer aux tolérances des cotes de la plaque PT102.\par\textbf{Il vous faut.} Un pied à coulisse 150 mm (celui de la mallette PT102) · Un coffret de cales (10, 40, 50, 60, 90 et 100 mm) · Une plaque PT102 (cotes a = 105 et o = 34 mm) · Ellistat\end{tcolorbox}
\begin{tcolorbox}[pred]Le pied à coulisse est juste à 10 mm. L'est-il autant à 150 mm ? Écrivez l'erreur que vous attendez à 150 mm, en µm.\reponse\end{tcolorbox}
\titre{À faire}
\begin{enumerate}[label=\textcolor{cyan}{\textbf{\arabic*.}},leftmargin=15pt,itemsep=1.5pt]
\item Fermer les becs, nettoyer, afficher le zéro. Nettoyer les cales.
\item Mesurer chaque longueur du tableau \textbf{5 fois}, en ouvrant et en refermant les becs à chaque mesure. Pour 140 mm, empiler 100 + 40 ; pour 150 mm, empiler 100 + 50 (adhérence, comme en M10).
\item Calculer pour chaque longueur la moyenne, puis le biais en µm : moyenne moins valeur de la cale.
\item \tagE\ \menu{Gage R\&R}, onglet \menu{Linéarité}, bouton « Ajouter une Etude de linéarité ». Remplir : Nom ; Nombre de produits 7 (Ellistat propose 4) ; Nombre de mesures par produit 5 ; Valeur k 3. Laisser les autres cases vides. Valider.
\item \tagE\ Onglet \menu{Mesures} : un bloc par longueur. Saisir la VALEUR VRAIE (10, 50, 60, 90, 100, 140, 150) puis ses 5 mesures dans la colonne MESURE.
\item \tagE\ Onglet \menu{Résultat} : lire « Justesse » et « Pente » (colonnes Valeur, t, Risque alpha, Significatif ?). Puis l'onglet \menu{Graphiques}, qui trace l'erreur selon la valeur vraie.
\end{enumerate}
\vspace{3pt}{\renewcommand{\arraystretch}{1.5}\small\setlength{\tabcolsep}{3pt}\begin{tabularx}{\linewidth}{|l|*{5}{>{\centering\arraybackslash}X|}>{\centering\arraybackslash}p{15mm}|>{\centering\arraybackslash}p{15mm}|}\hline
\textbf{Longueur} & \scriptsize 1 & \scriptsize 2 & \scriptsize 3 & \scriptsize 4 & \scriptsize 5 & \textbf{\scriptsize Moyenne} & \textbf{\scriptsize Biais (µm)}\\\hline
10 mm &  &  &  &  &  &  & \\\hline
50 mm &  &  &  &  &  &  & \\\hline
60 mm &  &  &  &  &  &  & \\\hline
90 mm &  &  &  &  &  &  & \\\hline
100 mm &  &  &  &  &  &  & \\\hline
140 mm &  &  &  &  &  &  & \\\hline
150 mm &  &  &  &  &  &  & \\\hline
\end{tabularx}}\par{\small Biais = moyenne $-$ valeur de la cale. 1 µm = 0,001 mm.}\par
\newpage
\titre{J'observe, je décide}
\begin{enumerate}[label=\textcolor{cyan}{\textbf{Q\arabic*.}},leftmargin=22pt]
\item Le biais est-il le même pour les sept longueurs ? Dans quel sens évolue-t-il avec la longueur ?\reponse[2]
\item Pente = \champ{18mm} mm par mm. Que dit Ellistat dans la colonne « Significatif ? » ? Plus grand biais relevé : \champ{14mm} µm.\reponse[1]
\item La cote $a = 105$ mm de la plaque a pour tolérance générale ISO 2768-mK, soit $\pm 0{,}3$ mm ($T = 0{,}6$ mm). Quelle part de $T$ représente le plus grand biais ? Faut-il s'en soucier ?\reponse[2]
\item La cote $o = 34 \pm 0{,}1$ mm ($T = 0{,}2$ mm). Appliquez la règle de M1 (résolution au plus 5 \% de $T$) : ce pied à coulisse convient-il ?\reponse[2]
\end{enumerate}
\begin{tcolorbox}[retenir]Le biais d'un pied à coulisse peut grandir avec la longueur. On le qualifie avec des étalons de la longueur où il sert, et on compare toujours à la tolérance de la cote.\end{tcolorbox}
\vfill{\footnotesize\color{uimm}\textbf{Ressources de la base :} Données de secours : \href{https://sjaubert.github.io/000_OI/SPC/Pages/Donnees_a_copier.html\#lin7}{\nolinkurl{https://sjaubert.github.io/000_OI/SPC/Pages/Donnees_a_copier.html}}}
\end{document}